% =====================================================================
%  preamble.tex — pengaturan tata letak skripsi (format FEB UNDIP)
%  Anda jarang perlu mengubah berkas ini. Edit data diri di frontmatter.tex
% =====================================================================

\usepackage{setspace}
\usepackage{fancyhdr}
\usepackage{graphicx}
\usepackage{float}
\usepackage{booktabs}
\usepackage{longtable}
\usepackage{caption}

% --- Penomoran bab gaya Indonesia: "BAB I", "BAB II", ... -------------
\renewcommand{\thechapter}{\Roman{chapter}}
\renewcommand{\chaptername}{BAB}

% KOMA-Script: tampilkan "BAB I" lalu judul bab di baris berikutnya,
% rata tengah, huruf kapital.
\renewcommand*{\raggedchapter}{\centering}
% Semua judul (bab, sub-bab) & entri daftar isi pakai huruf serif (Times),
% bukan sans-serif bawaan KOMA-Script.
\addtokomafont{disposition}{\rmfamily}
\setkomafont{chapter}{\normalfont\bfseries\large\MakeUppercase}
\renewcommand*{\chapterformat}{\chapapp~\thechapter\\[0.3em]}
\renewcommand*{\chapterlinesformat}[3]{#2#3}

% --- Nomor tabel & gambar mengikuti nomor bab (arabik): 4.1, 4.2 ------
\renewcommand{\thetable}{\arabic{chapter}.\arabic{table}}
\renewcommand{\thefigure}{\arabic{chapter}.\arabic{figure}}
\renewcommand{\theequation}{\arabic{chapter}.\arabic{equation}}

% --- Header / footer --------------------------------------------------
\pagestyle{fancy}
\fancyhf{}
\fancyhead[R]{\thepage}
\renewcommand{\headrulewidth}{0pt}

% --- Keterangan tabel/gambar ------------------------------------------
\captionsetup{font=normalsize, labelsep=space}

% --- Matikan halaman judul bawaan Quarto (kita pakai sampul sendiri) ---
\AtBeginDocument{\let\maketitle\relax}

% Spasi 1,5 sudah diatur lewat linestretch di _quarto.yml.
% Indentasi awal paragraf:
\setlength{\parindent}{1.27cm}
